% ---------------------------------------------------
% Trabajo Final de Grado
% Author: Fabián González Lence <alu0101549491@ull.edu.es>
% Chapter: Resultados
% ----------------------------------------------------

\chapter{Resultados} \label{chap:Results}

Los cinco proyectos desarrollados en el experimento están disponibles \cite{URL::TFG-Web} de forma pública y accesible: los tres primeros exclusivamente siendo un \textit{frontend} estático, y los dos últimos incluyendo \textit{backend} desplegado en la nube. Este capítulo documenta la infraestructura de despliegue adoptada (Sección~\ref{sec:results-deployment}), las métricas de calidad de código obtenidas mediante SonarQube y el análisis del agente de \ac{IA} (Sección~\ref{sec:results-code-quality}), y una valoración global de los proyectos resultantes (Sección~\ref{sec:results-analysis}).

\section{Despliegue de los proyectos} \label{sec:results-deployment}

Los cinco proyectos se despliegan desde el mismo repositorio monorepo bajo dos estrategias según su naturaleza. Los cinco \textit{frontends} se publican como sitios estáticos en GitHub Pages mediante un flujo de GitHub Actions que construye cada proyecto en secuencia, inyecta la variable \texttt{BASE\_URL} en el fichero \texttt{vite.config.ts} de cada uno para que se sirva bajo su propio subdirectorio, y se genera una página de entrada que enlaza a los cinco subproyectos \cite{URL::TFG-Web}. Los \textit{backends} de \ac{CARTO} y \ac{TENNIS} se desplegaron en Render\footnote{\href{https://render.com}{Render -- \texttt{render.com}}}, con la base de datos alojada en Supabase\footnote{\href{https://supabase.com}{Supabase -- \texttt{supabase.com}}}. Para \ac{CARTO} se exploró inicialmente Railway\footnote{\href{https://railway.app}{Railway -- \texttt{railway.app}}}, pero las restricciones de su plan gratuito y la incompatibilidad con la estructura monorepo motivaron el cambio. Ambos servicios se definen en el fichero \texttt{render.yaml} de la raíz del repositorio, centralizando la gestión de infraestructura en el mismo. 
Durante la migración de \ac{CARTO} a Supabase, las credenciales de conexión de la base de datos fueron expuestas accidentalmente en el repositorio público a causa de un error cometido por el agente de GitHub Copilot, esa exposición fue detectada por GitGuardian\footnote{\href{https://www.gitguardian.com}{GitGuardian -- \texttt{gitguardian.com}}} y las claves fueron invalidadas y regeneradas de inmediato.

\section{Métricas de calidad del código} \label{sec:results-code-quality}

La calidad del código se evaluó mediante dos herramientas complementarias. La primera es SonarQube Cloud\footnote{\href{https://sonarcloud.io/organizations/tfg-fabian-gonzalez-lence}{Organización del \texttt{TFG} en SonarQube Cloud\\ \hspace*{1.9em}\texttt{sonarcloud.io/organizations/tfg-fabian-gonzalez-lence}}}, integrado en el repositorio vía GitHub Actions con un fichero \texttt{sonar-project.properties} por proyecto. La segunda es un agente personalizado de GitHub Copilot denominado \textit{Code Quality Agent}, generado mediante \textit{meta-prompting} con Claude Opus 4.7, debido a su gran ventana de contexto. 
El agente aplica el estándar ISO/IEC 25010 \cite{ISO25010:2023} recorriendo las mismas métricas de calidad que SonarCloud (fiabilidad, seguridad, mantenibilidad, etc.), con el objetivo de contrastar ambas evaluaciones. Las Tablas~\ref{tab:sonarqube-issues} y~\ref{tab:sonarqube-ratings} recogen las métricas obtenidas directamente de la \ac{API} de SonarCloud, mientras que la Tabla~\ref{tab:quality-results} recoge el análisis complementario realizado por el agente de \ac{IA} según el estándar ISO/IEC 25010. La cobertura de \ac{CARTO} y \ac{TENNIS} figura como no disponible porque ambos proyectos utilizan exclusivamente pruebas \textit{end-to-end} con Playwright, cuya cobertura no es compatible con el formato LCOV de SonarCloud.

\begin{table}[H]
\centering
\small
\renewcommand{\arraystretch}{1.2}
\begin{tabular}{|>{\raggedright\arraybackslash}p{2cm}|>{\centering\arraybackslash}p{1.4cm}|>{\centering\arraybackslash}p{3.4cm}|>{\centering\arraybackslash}p{2.4cm}|>{\centering\arraybackslash}p{2cm}|}
\hline
\textbf{Proyecto} & \textbf{Bugs} & \textbf{Vulnerabilidades} & \textbf{\textit{Code smells}} & \textbf{Cobertura} \\
\hline
\ac{HANGMAN} & 0  & 0 & 20  & 95,8\,\% \\
\hline
\ac{PLAYER}  & 0  & 0 & 81  & 83,3\,\% \\
\hline
\ac{BALATRO} & 4  & 2 & 91  & 72,0\,\% \\
\hline
\ac{CARTO}   & 7  & 8 & 440 & ---      \\
\hline
\ac{TENNIS}  & 61 & 0 & 561 & ---      \\
\hline
\end{tabular}
\caption{Incidencias y cobertura de tests según SonarCloud por proyecto.}
\label{tab:sonarqube-issues}
\end{table}

\begin{table}[H]
\centering
\small
\renewcommand{\arraystretch}{1.2}
\begin{tabular}{|>{\raggedright\arraybackslash}p{1.8cm}|>{\centering\arraybackslash}p{2cm}|>{\centering\arraybackslash}p{2cm}|>{\centering\arraybackslash}p{3.2cm}|>{\centering\arraybackslash}p{2.6cm}|}
\hline
\textbf{Proyecto} & \textbf{Fiabilidad} & \textbf{Seguridad} & \textbf{Mantenibilidad} & \textbf{\textit{Quality Gate}} \\
\hline
\ac{HANGMAN} & A & A & A & \texttt{OK} \\
\hline
\ac{PLAYER}  & A & A & A & \texttt{OK} \\
\hline
\ac{BALATRO} & B & B & A & \texttt{OK} \\
\hline
\ac{CARTO}   & C & C & A & \texttt{OK} \\
\hline
\ac{TENNIS}  & E & A & A & \texttt{ERROR} \\
\hline
\end{tabular}
\caption{Calificaciones de calidad según SonarCloud. Fiabilidad, Seguridad y Mantenibilidad en una escala de A--E, donde A es la mejor calificación y E la peor.}
\label{tab:sonarqube-ratings}
\end{table}

\begin{table}[H]
\centering
\small
\renewcommand{\arraystretch}{1.2}
\begin{tabular}{|>{\raggedright\arraybackslash}p{1.8cm}|>{\centering\arraybackslash}p{2.1cm}|>{\centering\arraybackslash}p{1.4cm}|>{\centering\arraybackslash}p{1.4cm}|>{\centering\arraybackslash}p{1.4cm}|>{\centering\arraybackslash}p{1.4cm}|>{\centering\arraybackslash}p{2cm}|}
\hline
\textbf{Proyecto} & \textbf{Bloqueante} & \textbf{Crítica} & \textbf{Mayor} & \textbf{Menor} & \textbf{Total} & \textbf{Valoración} \\
\hline
\ac{HANGMAN} & 0 & 0 & 1  & 5  & 6   & Bueno \\
\hline
\ac{PLAYER}  & 1 & 2 & 9  & 12 & 24  & Mejorable \\
\hline
\ac{BALATRO} & 0 & 3 & 9  & 16 & 28  & Mejorable \\
\hline
\ac{CARTO}   & 5 & 6 & 13 & 6  & 30  & Crítico \\
\hline
\ac{TENNIS}  & 0 & 5 & 12 & 8  & 25  & Crítico \\
\hline
\textbf{Total} & \textbf{6} & \textbf{16} & \textbf{44} & \textbf{47} & \textbf{113} & --- \\
\hline
\end{tabular}
\caption{Recuento de incidencias y valoración según el análisis del agente de \ac{IA}.}
\label{tab:quality-results}
\end{table}

Ambas herramientas convergen en la tendencia principal: la calidad decrece a medida que aumenta la complejidad del proyecto. El resultado más destacable en positivo es que la mantenibilidad obtiene calificación A en los cinco proyectos, coherente con la aplicación sistemática de los principios \ac{SOLID} y la \textit{Guía de Estilo de Google para TypeScript} en todas las fases del flujo. La cobertura describe una curva descendente en los tres proyectos con \textit{testing} unitario ---95{,}8\,\%, 83{,}3\,\% y 72{,}0\,\%--- que refleja la complejidad creciente del dominio.

La divergencia entre las dos herramientas es especialmente reveladora en los dos proyectos \textit{fullstack}. \ac{TENNIS} acumula 61 bugs y es el único que no supera el \textit{Quality Gate} ---fiabilidad E---, mientras que el agente de \ac{IA} le asigna cero incidencias bloqueantes frente a las cinco que detecta en \ac{CARTO}. La explicación reside en la naturaleza de los defectos: los 61 bugs de \ac{TENNIS} son errores de lógica de negocio en los estados de los partidos; las cinco incidencias bloqueantes de \ac{CARTO} son problemas de diseño arquitectónico cuya gravedad requiere una comprensión global del sistema que el análisis estático no captura con la misma precisión.

La progresión de \textit{code smells} revela que la deuda de mantenibilidad no escala linealmente: entre \ac{BALATRO} y \ac{CARTO} casi se quintuplican mientras los ficheros se multiplican por $\approx$2{,}5, lo que apunta a que la arquitectura \textit{fullstack} con \textit{Clean Architecture} genera deuda a un ritmo superior al crecimiento de la base de código. Asimismo, la revisión de \ac{TENNIS} resolvió 58 de las 60 incidencias documentadas, pero los 61 bugs del módulo de partidos no afloraron hasta la ejecución de las pruebas \textit{end-to-end}, evidencia de que el análisis estático no basta para garantizar la fiabilidad en flujos de negocio con estados complejos.

\section{Análisis de los proyectos resultantes} \label{sec:results-analysis}

Los cinco proyectos cumplen funcionalmente con las especificaciones de requisitos definidas en el Capítulo~\ref{chap:Design} y están accesibles en producción \cite{URL::TFG-Web}. La degradación de calidad documentada en la Sección~\ref{sec:results-code-quality} es coherente con la literatura sobre generación de código con \ac{LLMs} revisada en el Capítulo~\ref{chap:RelatedWork}: la calidad decrece a medida que crece el tamaño del contexto necesario, y la progresión de escala entre proyectos reproduce ese patrón con precisión. Las incidencias de mantenibilidad —complejidad ciclomática, funciones extensas, uso recurrente de código cuestionable como el \texttt{as any}— son sistemáticas y predecibles, mientras que las de seguridad y fiabilidad —almacenamiento inseguro de credenciales, ausencia de límites de error en los árboles de componentes— solo emergen de una comprensión global del diseño del sistema, algo que los modelos actuales gestionan de forma inconsistente. Las implicaciones de este comportamiento se desarrollan en el Capítulo~\ref{chap:Conclusiones}.